\textbf{Problème 1.7 (Dujella, 2026).} Soit $D_{m,n}(N)$ le nombre de $D(n)$-$m$-uplets à éléments dans l'ensemble $\{1,2,\ldots,N\}$, pour un entier positif $N$ donné.
\begin{enumerate}
\item Si $n$ est un carré parfait, existe-t-il une constante $C(n)$ ne dépendant que de $n$ telle que $D_{3,n} \sim C(n)N \log N$ ? Est-il vrai que $C(n) = C(1) = \frac{3}{\pi^2}$ ?
\item Si $n$ n'est pas un carré parfait, existe-t-il une constante $C(n)$ ne dépendant que de $n$ telle que $D_{3,n} \sim C(n)N$ ?
\item Si $n$ est un carré parfait, que peut-on dire du comportement asymptotique de $D_{4,n}$ ?
\item Si $n$ n'est pas un carré parfait, est-il vrai que $D_{4,n}$ est fini ?
\end{enumerate}

\medskip

\textit{Remarque.} Une réponse affirmative à la question (b), dans le cas où $|n|$ est premier ou $n=-1$, est donnée dans Adžaga, N., Dražić, G., Dujella, A., Pethő, A. : Asymptotics of $D(q)$-pairs and triples via $L$-functions of Dirichlet characters. Ramanujan J. \textbf{66}, Article 11 (2025).

Les questions (a) et (b) sont entièrement résolues dans Badesa, J. : On the asymptotics of $D(n)$-pairs and triples. Ramanujan J. \textbf{67}, Article 101 (2025) ; les réponses aux deux questions sont affirmatives. En particulier, pour $n$ carré parfait, $D_{3,n} \sim \frac{3}{\pi^2}N \log N$.

Une preuve alternative de (a) et (b) est donnée dans l'article récent Dražić, G., Kazalicki, M., Mrazović, R. : Equidistribution of Diophantine pairs among the equivalence classes of quadratic forms. Prépublication (2025). arXiv : 2512.22902.
